\chapter{Related Work}
\label{chap:related_work}

\section{Full-Band Speech Enhancement}
\label{sec:related_observed_band_gap}

Neural speech enhancement (SE) has evolved from time-frequency mask estimation to phase-aware, complex-domain, time-domain, and modern TF-domain modeling. Early systems commonly estimated real-valued masks in the STFT domain, while later methods introduced phase-sensitive masks, complex ratio masks, and complex-valued networks to better handle phase-related distortion. Time-domain systems replaced the fixed STFT representation with learned analysis and synthesis bases, and recent TF-domain models further improved enhancement quality by explicitly modeling temporal and spectral dependencies~\cite{erdogan2015psm,williamson2016cirm,hu2020dccrn,luo2019convtasnet,cao2022cmgan,lu2023mpsenet,saijo2024tflocoformer}. However, many reported comparisons remain tied to 16 kHz or protocol-specific settings, where the upper-band components present in 48 kHz speech are outside the input representation.

The motivation for full-band SE is also supported by perceptual studies on high-frequency speech information. Extended high-frequency cues above the conventional wideband range have been shown to contribute to speech perception in noise and to the perception of fricative-related cues~\cite{stelmachowicz2001,monson2014,moore2016}. These findings motivate retaining and refining upper-band observations when the full-band spectrum is already present in the input.

Speech bandwidth extension and speech super-resolution address a different information condition from direct full-band SE. These tasks aim to reconstruct or supplement missing high-frequency content from band-limited inputs, often targeting a higher output sampling rate or wider output bandwidth~\cite{liu2022voicefixer,andreev2022hifiplusplus,audiosr}. In contrast, 48 kHz full-band SE starts from a noisy full-band observation, where the upper-band spectrum is already present but corrupted by noise. Therefore, full-band SE requires noise suppression and spectral refinement over the observed full-band signal, rather than only high-band reconstruction from a band-limited input. This observed-band condition motivates architectures that retain noisy full-band evidence while using reliable low-band speech structure as refinement context.

\section{Efficient Full-Band Modeling}
\label{sec:related_practical_frequency_structured}

Practical full-band SE has been studied mainly under real-time and low-complexity constraints. PercepNet enhances 48 kHz speech using perceptually motivated critical-band features and comb filtering, while DeepFilterNet2 combines ERB-domain enhancement with deep filtering for low-complexity full-band enhancement~\cite{valin2020percepnet,schroter2022deepfilternet2}. Such perceptual or ERB-band representations trade frequency resolution for efficiency, motivating complementary TF-domain designs that retain a closer connection to the observed STFT while still controlling computational cost.

Frequency-structured modeling has also become a major design axis in modern SE and source separation. Subband and band-split models process spectra through grouped frequency regions, while TF-domain dual-path models separate temporal and spectral modeling to improve efficiency and representation power. BSRNN and BS-RoFormer show that band-wise modeling is effective for high-resolution audio modeling, and TF-Locoformer develops an RNN-free TF-domain dual-path architecture with local modeling by convolution~\cite{luo2022bsrnn,lu2023bsroformer,saijo2024tflocoformer}. Lightweight systems such as GTCRN, LiSenNet, and FastEnhancer further show that compact recurrent, subband, and frequency-time modeling can improve the performance-efficiency trade-off in real-time SE~\cite{gtcrn,lisennet,ahn2025fastenhancer}.

Many band-split and lightweight models emphasize compact subband grouping or efficient within-band/global sequence modeling. TF-Rehancer also uses frequency compression: its input stem applies a learned frequency-strided embedding to reduce token count. The difference is that the final enhancement remains tied to the original noisy STFT through full-resolution complex filtering. This design targets direct 48 kHz full-band SE, where the high-frequency region is not missing but available as noisy spectral evidence in the decoder path.

\section{Query Refinement and TF Filtering}
\label{sec:related_query_filtering}

Query-based restoration provides another relevant direction. TF-Restormer introduces an asymmetric encoder-decoder architecture for speech restoration with decoupled input and output sampling rates. Its encoder focuses on the observed input bandwidth, while the decoder restores target frequency regions through frequency extension queries~\cite{shin2026tfre}. This mechanism is well matched to bandwidth restoration, where decoder-side queries can serve as slots for frequency regions that are absent from the input.

Direct 48 kHz full-band SE has a different information condition. TF-Restormer uses decoder-side queries for missing or target frequency regions in restoration, whereas TF-Rehancer uses query modeling in direct full-band SE, where the high-band region is observed and noisy. The query is therefore not a missing-band placeholder; it is derived from observed high-band evidence and updated using low-band context.

Filtering-based enhancement is closely related to this observation-preserving view. Instead of directly mapping a latent representation to a clean spectrum, deep filtering estimates complex local filters and applies them to the noisy STFT~\cite{schroter2022deepfilternet2}. Recent correlation-to-filter approaches such as SR-CorrNet provide related evidence that correlation-derived or structure-conditioned filters can be effective alternatives to direct spectral mapping~\cite{shin2026srcorrnet}. TF-Rehancer follows this filtering-based principle by estimating the enhanced spectrum through complex filtering over the original noisy STFT.

Prior work covers complementary parts of this problem. Practical full-band systems emphasize perceptual-band efficiency, frequency-structured models mainly address efficient within-band or band-split modeling, and query-based restoration methods mainly target missing-band reconstruction. In contrast, TF-Rehancer targets direct full-band SE through low-band-guided full-band refinement and complex filtering over the noisy STFT.
